%% Allgemeines -----------------------------------

% Fortlaufende Nummerierung - Entweder 1, 2, 3, 4, ... oder mit einem semantischen Schema
\def \decisionID {4}

% Datum, an dem der Record anefertigt wurde
\def \decisionDate {22.11.2026}

% Wer hat an dem Record mitgewirkt?
\def \decisionParticipants {Niklas Käsehagen, OneTutor (Customer Proxy)}

% Der Status in dem sich die Anforderung in diesem Record aktuell befindet.
% Möglichkeiten: Vorgeschlagen, Akzeptiert, Abgelehnt, Abgelöst
\def \decisionStatus {Abgelöst}

% ID des Records, der diesen Record ablöst, falls zutreffend
\def \decisionSuperseded {6, 7}

% ID des Records, den dieser Record ablöst, falls zutreffend
\def \decisionSuperseding {3}

%% Anforderung -----------------------------------

% Titel der Anforderung
\def \requirementTitle {Ähnlichkeitsmaß muss entscheiden können, welche Folien ähnlich zueinander sind}

% Kontext der Anforderung
\def \requirementContext {Wenn ein Foliensatz in das Repository hochgeladen wird, soll die Software Similar Slides erkennen können,
ob die Folien jeweils bereits im Repository enthalten sind oder ob sich ähnliche Folien im Repository befinden. Kann dies nicht unterschieden werden, so unterscheidet sich das Repository nicht von der aktuellen Speicherung der Foliensätze und die Dozierenden müssen viel Zeit aufwenden, um Folien zu suchen und zu organisieren.}

% Konsequenz der Anforderung
\def \requirementConsequences {Das System prüft die mögliche Ähnlichkeit jeder Folie mit Folien im Repository. Dabei liegt die Priorität zunächst beim Text der Folien, danach beim Titel der Folien und zuletzt bei möglichen Graphiken und Bildern. Ab welchem Ähnlichkeitsmaß zwei Folien als ähnlich gelten ist nicht fest vorgegeben und soll während der Entwicklung erarbeitet werden.}
